% DERIVED from formal_layer.json — read-only, do not edit.
\begin{assumptionv}{ass:uniform-covariates}[Uniform participant strata]
The participant stratum \(X\) is uniformly distributed on \([d]\):
\[
P(X=j)=\frac{1}{d}\qquad(j\in[d]).
\]
\end{assumptionv}

\begin{assumptionv}{ass:fair-randomization}[Fair conditional treatment assignment]
The binary treatment assignment \(A\), conditional on the participant stratum \(X\) and potential outcomes \(Y^0,Y^1\), satisfies
\[
\mathcal L_P(A\mid X,Y^0,Y^1)=\operatorname{Bernoulli}(1/2).
\]
\end{assumptionv}

\begin{assumptionv}{ass:consistency}[Observed outcome consistency]
The observed outcome \(Y\) equals \(Y^{\mathrm{pot}}_A\), the potential outcome selected by the realized assignment, almost surely:
\[
P(Y=Y^{\mathrm{pot}}_A)=1.
\]
\end{assumptionv}

\begin{assumptionv}{ass:interior-means}[Interior potential-outcome means]
The conditional potential-outcome means \(\mu_{aj}(P)\) satisfy
\[
\mu_{aj}(P)\in[1/4,3/4]
\qquad(a\in\{0,1\},\ j\in[d]).
\]
\end{assumptionv}

\begin{assumptionv}{ass:iid-people}[Independent identically distributed participants]
The vector of observed participant records \(\mathbf O\) has the product law
\[
\mathcal L_P(\mathbf O)=P_O^{\otimes n},
\]
where \(P_O\) is the observed-record marginal law induced by \(P\).
\end{assumptionv}

\begin{definitionv}{synth_1}[Full-data causal law space]
For an integer \(d\ge2\), write \([d]=\{1,\ldots,d\}\). Define \(\mathcal P_d\) as the class of all probability laws on the finite record space
\[
[d]\times\{0,1\}^4,
\]
equipped with its discrete sigma-algebra. The coordinate variables are ordered as \((X,A,Y,Y^0,Y^1)\): \(X\in[d]\) is the covariate stratum, \(A\in\{0,1\}\) is the treatment assignment, \(Y\in\{0,1\}\) is the observed outcome, and \(Y^0,Y^1\in\{0,1\}\) are the potential outcomes under control and treatment, respectively. Write \(Y^{\mathrm{pot}}_0=Y^0\) and \(Y^{\mathrm{pot}}_1=Y^1\). For \(P\in\mathcal P_d\), the observed record is \(O=(X,A,Y)\in\mathcal O=[d]\times\{0,1\}^2\), and \(P_O\) denotes its marginal law under \(P\).
\end{definitionv}

\begin{definitionv}{def:causal-model}[Causal law class $\mathcal V_d$]
\(\mathcal V_d\) is the class of full-data probability laws \(P\in\mathcal P_d\) satisfying all four conditions:
\begin{itemize}
\item \textbf{(Uniform strata.)} \Cref{obj:ass:uniform-covariates}.
\item \textbf{(Fair assignment.)} \Cref{obj:ass:fair-randomization}.
\item \textbf{(Consistency.)} \Cref{obj:ass:consistency}.
\item \textbf{(Interior means.)} \Cref{obj:ass:interior-means}.
\end{itemize}
\end{definitionv}

\begin{assumptionv}{ass:independent-randomness}[Independent public and analyst randomness]
The shared public seed \(R\), with probability law \(\lambda\), and the analyst randomizer \(U\) have conditional joint law
\[
\mathcal L(R,U\mid\mathbf O)
=\lambda\otimes\operatorname{Uniform}[0,1],
\]
where \(\mathbf O\) is the vector of observed participant records.
\end{assumptionv}

\begin{assumptionv}{ass:full-record-privacy}[Full-record local privacy]
For every participant index \(i\), observed records \(o,o'\), history \(\eta\), public seed \(r\), and measurable message set \(E\) in their declared domains,
\[
Q_i(E\mid o,\eta,r)\le e^\varepsilon Q_i(E\mid o',\eta,r).
\]
\end{assumptionv}

\begin{definitionv}{synth_2}[Participant message histories]
For a protocol with \(n\) participants, let \((\mathcal Z_i)_{i=1}^n\) be its standard-Borel participant message spaces. For each \(i\in\{1,\ldots,n\}\), define the history space
\[
\mathcal H_i=\prod_{\ell=1}^{i-1}\mathcal Z_\ell
\]
with the product sigma-algebra; the empty product \(\mathcal H_1\) is a singleton. A history \(\eta\in\mathcal H_i\) records the messages of participants \(1,\ldots,i-1\). Thus, with public-seed space \(\mathcal R\), the participant kernel \(Q_i\) is a Markov kernel from \(\mathcal O\times\mathcal H_i\times\mathcal R\) to \(\mathcal Z_i\), and the realized history preceding participant \(i\) is \((Z_1,\ldots,Z_{i-1})\).
\end{definitionv}

\begin{assumptionv}{ass:noninteraction}[Independence from message history]
For every participant index \(i\), observed record \(o\), public seed \(r\), and histories \(\eta,\eta'\in\mathcal H_i\),
\[
Q_i(\cdot\mid o,\eta,r)=Q_i(\cdot\mid o,\eta',r).
\]
\end{assumptionv}

\begin{definitionv}{def:sequential-class}[Sequential protocol class $\mathfrak Q_{\mathrm{SI}}(n,d,\varepsilon)$]
\(\mathfrak Q_{\mathrm{SI}}(n,d,\varepsilon)\) is the class of sequences
\[
Q=(Q_i)_{i=1}^n
\]
of the declared Markov kernels satisfying \cref{obj:ass:full-record-privacy}.
\end{definitionv}

\begin{definitionv}{def:noninteractive-class}[Noninteractive protocol class $\mathfrak Q_{\mathrm{NI}}(n,d,\varepsilon)$]
\(\mathfrak Q_{\mathrm{NI}}(n,d,\varepsilon)\) is the subclass of \(\mathfrak Q_{\mathrm{SI}}(n,d,\varepsilon)\) satisfying both \cref{obj:ass:full-record-privacy,obj:ass:noninteraction}.
\end{definitionv}

\begin{definitionv}{def:risk}[Minimax welfare risk $\mathfrak R_{\mathcal C}(n,d,\varepsilon)$]
\[
\mathfrak R_{\mathcal C}(n,d,\varepsilon)
=
\inf_{Q\in\mathfrak Q_{\mathcal C}(n,d,\varepsilon)}
\inf_T
\sup_{P\in\mathcal V_d}
\mathbb E_{P,Q}
\left[
\bigl(T(\mathbf Z,R,U)-V(P)\bigr)^2
\right].
\]
Here \(\mathcal C\in\{\mathrm{NI},\mathrm{SI}\}\) selects the corresponding declared protocol class \(\mathfrak Q_{\mathcal C}\), and the second infimum ranges over Borel real-valued welfare estimates \(T(\mathbf Z,R,U)\) based on the private message vector \(\mathbf Z\), the public seed \(R\), and the independent uniform analyst randomizer \(U\). The target \(V(P)\) is the average of the larger arm mean in each stratum under \(P\).
\end{definitionv}

\begin{definitionv}{def:honest-length}[Honest interval length $\mathfrak H_{\mathcal C}(n,d,\varepsilon)$]
The minimax expected length of a globally honest welfare interval is
\[
\mathfrak H_{\mathcal C}(n,d,\varepsilon)
=
\inf_{Q\in\mathfrak Q_{\mathcal C}(n,d,\varepsilon)}
\ \inf_{\substack{I:\ \inf_{P\in\mathcal V_d}
\Pr_{P,Q}\{V(P)\in I(\mathbf Z,R,U)\}\ge0.90}}
\ \sup_{P\in\mathcal V_d}
\mathbb E_{P,Q}\!\left[
\operatorname{length} I(\mathbf Z,R,U)
\right].
\]
The second infimum ranges over connected closed intervals whose ordered endpoints are Borel functions of the private messages \(\mathbf Z\), public seed \(R\), and analyst randomizer \(U\), and take values in \([1/4,3/4]\).
\end{definitionv}

\begin{definitionv}{synth_3}[Private resource rates and their regimes]
For \(n,d\in\{2,3,\ldots\}\) and \(0<\varepsilon\le1\), put \(t=n\varepsilon^2\) and \(L=\log(ed)\). Define
\[
\rho(t,d)=
\begin{cases}
\dfrac{d^2}{tL},&t\ge d^2L,\\[4pt]
\displaystyle\min\left\{1,
\left[\dfrac{\log(e+d^2L/t)}{L}\right]^2\right\},
&0<t<d^2L,
\end{cases}
\qquad
r_{\mathcal C}(n,d,\varepsilon)=\rho(t,d),
\qquad
h_{\mathcal C}(n,d,\varepsilon)=\sqrt{\rho(t,d)}
\]
for each \(\mathcal C\in\{\mathrm{NI},\mathrm{SI}\}\).

For positive numerical sequences \(f_\ell,g_\ell\), write \(f_\ell\asymp g_\ell\) when there exist constants \(0<c\le C_0<\infty\) and an integer \(\ell_0\) such that
\[
c g_\ell\le f_\ell\le C_0g_\ell
\qquad(\ell\ge\ell_0).
\]
We say that these rates satisfy the resource conclusions when all the following properties hold.

For every integer \(d\ge2\) and every \(0<t<d^2\log(ed)\),
\[
\rho(t,d)=1
\quad\Longleftrightarrow\quad
t\le\frac{d^2\log(ed)}{ed-e}.
\]
Thus the saturated regime ends at \(d^2\log(ed)/(ed-e)\), the intermediate regime occupies
\[
\frac{d^2\log(ed)}{ed-e}<t<d^2\log(ed),
\]
and the dense regime starts at \(t=d^2\log(ed)\).

For every sequence \((n_\ell,d_\ell,\varepsilon_\ell)_{\ell\ge1}\) satisfying
\[
n_\ell,d_\ell\in\{2,3,\ldots\},
\qquad 0<\varepsilon_\ell\le1,
\]
put \(t_\ell=n_\ell\varepsilon_\ell^2\) and \(L_\ell=\log(ed_\ell)\). Then
\[
\rho(t_\ell,d_\ell)\longrightarrow0
\quad\Longleftrightarrow\quad
t_\ell\longrightarrow\infty
\ \text{and}\
\frac{\log(1+d_\ell^2/t_\ell)}{L_\ell}\longrightarrow0,
\]
and
\[
\rho(t_\ell,d_\ell)\asymp t_\ell^{-1}
\quad\Longleftrightarrow\quad
\left(\exists a>0,\ \exists\ell_0,\ \forall\ell\ge\ell_0:
t_\ell\ge a\right)
\ \text{and}\
\left(\exists M<\infty,\ \exists\ell_1,\ \forall\ell\ge\ell_1:
\min\{t_\ell,d_\ell\}\le M\right).
\]
For every such sequence with \(t_\ell\to\infty\),
\[
\rho(t_\ell,d_\ell)\asymp t_\ell^{-1}
\quad\Longleftrightarrow\quad
\exists M\in\mathbb N,\ \exists\ell_0,\ \forall\ell\ge\ell_0:
d_\ell\le M.
\]

Finally, for every \(0<a\le b<\infty\), there exist constants \(0<c\le C_0<\infty\) and an integer \(D_0\) such that, for every integer \(d\ge\max\{2,D_0\}\) and every \(t>0\) with \(a\le t/d^2\le b\),
\[
c\left[\frac{\log\log(ed)}{\log(ed)}\right]^2
\le \rho(t,d)\le
C_0\left[\frac{\log\log(ed)}{\log(ed)}\right]^2.
\]
The constants and dimension threshold in this last comparison may depend on \(a,b\).
\end{definitionv}

\begin{definitionv}{def:symmetric-law}[Symmetric causal law $G_\theta$]
The full-data law \(G_\theta\) is generated by
\[
X\sim\operatorname{Uniform}[d],
\qquad
Y^0\mid X=j\sim\operatorname{Bernoulli}\!\left(\frac{1-\theta_j}{2}\right),
\qquad
Y^1\mid X=j\sim\operatorname{Bernoulli}\!\left(\frac{1+\theta_j}{2}\right),
\]
with \(Y^0\) and \(Y^1\) conditionally independent given \(X\). Independently of \((X,Y^0,Y^1)\), generate the treatment assignment and set the observed outcome by
\[
A\sim\operatorname{Bernoulli}(1/2),
\qquad
Y=Y^{\mathrm{pot}}_A.
\]
\end{definitionv}

\begin{definitionv}{synth_6}[Sequence conclusions for paired minimax values]
For each \(\mathcal C\in\{\mathrm{NI},\mathrm{SI}\}\), let
\[
\mathcal R_{\mathcal C}(n,d,\varepsilon)
=\mathfrak R^{\mathrm{TV}}_{\mathcal C}
(\mathcal P_d^{\mathrm{pair}};n,d,\varepsilon),
\qquad
\mathcal H_{\mathcal C}(n,d,\varepsilon)
=\mathfrak H^{\mathrm{TV}}_{\mathcal C}
(\mathcal P_d^{\mathrm{pair}};n,d,\varepsilon),
\]
where the minimax risk and globally \(0.90\)-honest connected-interval length are defined in \cref{obj:def:tv-private-decision-values}. We say that the sequence-quantified value conclusions hold when the following properties hold for each class.

For every sequence \((n_\ell,d_\ell,\varepsilon_\ell)_{\ell\ge1}\) with
\[
n_\ell,d_\ell\in\{2,3,\ldots\},
\qquad 0<\varepsilon_\ell\le1,
\]
write
\[
t_\ell=n_\ell\varepsilon_\ell^2,\qquad
L_\ell=\log(ed_\ell),\qquad
R_\ell=\mathcal R_{\mathcal C}(n_\ell,d_\ell,\varepsilon_\ell),
\qquad
H_\ell=\mathcal H_{\mathcal C}(n_\ell,d_\ell,\varepsilon_\ell).
\]
Then
\[
R_\ell\longrightarrow0
\quad\Longleftrightarrow\quad
t_\ell\longrightarrow\infty
\ \text{and}\
\frac{\log(1+d_\ell^2/t_\ell)}{L_\ell}\longrightarrow0,
\]
and, separately,
\[
H_\ell\longrightarrow0
\quad\Longleftrightarrow\quad
t_\ell\longrightarrow\infty
\ \text{and}\
\frac{\log(1+d_\ell^2/t_\ell)}{L_\ell}\longrightarrow0.
\]

Define \(R_\ell\asymp t_\ell^{-1}\) to mean that there exist \(0<c\le C_0<\infty\) and \(\ell_0\) such that
\[
c/t_\ell\le R_\ell\le C_0/t_\ell
\qquad(\ell\ge\ell_0).
\]
For every allowed sequence,
\[
R_\ell\asymp t_\ell^{-1}
\quad\Longleftrightarrow\quad
\left(\exists a>0,\ \exists\ell_0,\ \forall\ell\ge\ell_0:
t_\ell\ge a\right)
\ \text{and}\
\left(\exists M<\infty,\ \exists\ell_1,\ \forall\ell\ge\ell_1:
\min\{t_\ell,d_\ell\}\le M\right).
\]
For every allowed sequence satisfying \(t_\ell\to\infty\),
\[
R_\ell\asymp t_\ell^{-1}
\quad\Longleftrightarrow\quad
\exists M\in\mathbb N,\ \exists\ell_0,\ \forall\ell\ge\ell_0:
d_\ell\le M.
\]

Finally, for every \(0<a\le b<\infty\), there exist \(0<c\le C_0<\infty\) and an integer \(D_0\) such that, for every allowed tuple \((n,d,\varepsilon)\) with \(d\ge D_0\) and
\[
a\le\frac{n\varepsilon^2}{d^2}\le b,
\]
both comparisons hold:
\[
c\left[\frac{\log\log(ed)}{\log(ed)}\right]^2
\le\mathcal R_{\mathcal C}(n,d,\varepsilon)
\le C_0\left[\frac{\log\log(ed)}{\log(ed)}\right]^2,
\]
\[
\sqrt c\left|\frac{\log\log(ed)}{\log(ed)}\right|
\le\mathcal H_{\mathcal C}(n,d,\varepsilon)
\le\sqrt{C_0}\left|\frac{\log\log(ed)}{\log(ed)}\right|.
\]
These constants and the dimension threshold may depend on \(a,b\) and the fixed class, and are uniform over the displayed resource domain.
\end{definitionv}

\begin{definitionv}{synth_5}[Resource-selected paired polynomial procedure]
Fix \(n,d\in\{2,3,\ldots\}\), \(0<\varepsilon\le1\), and \(C_0>0\). Put
\[
t=n\varepsilon^2,\qquad L=\log(ed),\qquad
\delta=\tanh(\varepsilon/2),\qquad b=d/\delta,
\qquad m=\lfloor n/3\rfloor.
\]
Use a singleton public seed. For participants \(i=1,\ldots,2m\), release a sign vector \(Z_i\in\{-1,1\}^d\) with conditional probabilities
\[
\mathsf K_i(z\mid(j,s))
=2^{-d}(1+\delta s z_j).
\]
The remaining participants release the constant vector \((-1,\ldots,-1)\). All releases use independent local randomness and are independent of histories.

If \(n=2\), set \(\widehat D=1/8\). If \(n>2\), use the first \(m\) active participants as the pilot block and the next \(m\) as the evaluation block. Define
\[
W^{\mathrm{pil}}_{ij}=bZ_{i,j},
\qquad
W^{\mathrm{ev}}_{ij}=bZ_{m+i,j},
\qquad 1\le i\le m,\quad j\in[d],
\]
and
\[
\sigma^2=\frac{b^2}{m},\qquad
\sigma=\sqrt{\sigma^2},\qquad
H=\log(e+\sigma^2L),\qquad
M_j=\frac1m\sum_{i=1}^mW^{\mathrm{pil}}_{ij},
\qquad
\overline W_j=\frac1m\sum_{i=1}^mW^{\mathrm{ev}}_{ij}.
\]
For the evaluation block, define the distinct-person moments
\[
\mathsf U_{0j}=1,\qquad
\mathsf U_{vj}
=\binom mv^{-1}
\sum_{\substack{J\subseteq\{1,\ldots,m\}\\ |J|=v}}
\prod_{i\in J}W^{\mathrm{ev}}_{ij},
\qquad 1\le v\le m.
\]
Their finite computation uses elementary-symmetric recursions.

For each even degree \(D\ge2\) selected below, define
\[
\mathcal T_0(x)=1,\qquad
\mathcal T_1(x)=x,\qquad
\mathcal T_{v+1}(x)=2x\mathcal T_v(x)-\mathcal T_{v-1}(x)
\quad(v\ge1),
\]
\[
p_D^{\mathrm{abs}}(x)
=\frac2\pi+\frac4\pi
\sum_{v=1}^{D/2}
\frac{(-1)^{v+1}\mathcal T_{2v}(x)}{4v^2-1}
=\sum_{v=0}^Dc_{D,v}x^v,
\]
and, for a positive radius \(a_{\mathrm{rad}}\),
\[
\widehat p_{a_{\mathrm{rad}},D,j}
=\sum_{v=0}^D
c_{D,v}a_{\mathrm{rad}}^{\,1-v}\mathsf U_{vj}.
\]

For \(n>2\), define \(\widehat F\) by the following ordered branches, taking the first applicable branch:
\[
\widehat F=
\begin{cases}
\displaystyle\frac1d\sum_{j=1}^d|\overline W_j|,
&L<4096,\\[6pt]
\displaystyle\frac1d\sum_{j=1}^d
\left[
\widehat p_{a_{\mathrm{rad}},D,j}
\mathbf1_{\{|M_j|\le T_*\}}
+\operatorname{sign}(M_j)\overline W_j
\mathbf1_{\{|M_j|>T_*\}}
\right],
&L\ge4096,\ t\ge d^2L,\\[6pt]
1/4,
&L\ge4096,\ t<d^2L,\ L/(1024H)<4,\\[4pt]
\displaystyle\frac1d\sum_{j=1}^d
\widehat p_{a_{\mathrm{rad}},D,j},
&L\ge4096,\ t<d^2L,\ L/(1024H)\ge4.
\end{cases}
\]
In the second branch, set
\[
D=2\left\lfloor\frac{L}{2048}\right\rfloor,
\qquad T_*=4\sigma\sqrt L,\qquad
a_{\mathrm{rad}}=8\sigma\sqrt L,
\qquad \operatorname{sign}(0)=0.
\]
In the fourth branch, set
\[
a_{\mathrm{rad}}=\frac12,
\qquad D=2\left\lfloor\frac{L}{2048H}\right\rfloor.
\]
For \(n>2\), set
\[
\widehat D=\max\left\{0,\min\left\{\frac14,\frac{\widehat F}{2}\right\}\right\}.
\]
For every \(n\ge2\), report
\[
q=\sqrt{10C_0\rho(t,d)},\qquad
J_{\mathrm{TV}}
=\left[
\max\{0,\widehat D-q\},
\min\{1/4,\widehat D+q\}
\right].
\]

We say that the concrete procedure guarantees hold with constant \(C_0\) when this fixed protocol belongs to \(\mathfrak K_{\mathrm{NI}}(n,d,\varepsilon)\) and, for every \(\theta\in[-1/2,1/2]^d\),
\[
\begin{aligned}
\mathbb E_{P_\theta,\mathsf K}
\left[(\widehat D-D_{\mathrm{TV}}(P_\theta))^2\right]
&\le C_0\rho(t,d),\\
\Pr_{P_\theta,\mathsf K}
\{D_{\mathrm{TV}}(P_\theta)\in J_{\mathrm{TV}}\}
&\ge0.90,\\
\mathbb E_{P_\theta,\mathsf K}
\left[\operatorname{length}J_{\mathrm{TV}}\right]
&\le C_0\sqrt{\rho(t,d)}.
\end{aligned}
\]
Here \(P_\theta(j,s)=(1+s\theta_j)/(2d)\) and
\(D_{\mathrm{TV}}(P_\theta)=(2d)^{-1}\sum_j|\theta_j|\).
Both decisions are Borel functions of the transcript, and the reported interval is nonempty, connected, closed, and contained in \([0,1/4]\).
\end{definitionv}

\begin{definitionv}{synth_4}[Finite sign-vector attainment]
Fix \(n,d\in\{2,3,\ldots\}\), \(0<\varepsilon\le1\), and \(C_0>0\). Let
\[
P_\theta(j,s)=\frac{1+s\theta_j}{2d},
\qquad
\theta\in[-1/2,1/2]^d,\quad j\in[d],\quad s\in\{-1,1\}.
\]
The paired family and its total-variation target are
\[
\mathcal P_d^{\mathrm{pair}}
=\{P_\theta:\theta\in[-1/2,1/2]^d\},
\qquad
D_{\mathrm{TV}}(P_\theta)
=\frac1{2d}\sum_{j=1}^d|\theta_j|.
\]
We say that the general attainment guarantees hold with constant \(C_0\) when there exist a protocol
\[
\mathsf K\in\mathfrak K_{\mathrm{NI}}(n,d,\varepsilon),
\]
a Borel estimator \(\widehat D\), and a Borel interval decision \(J_{\mathrm{TV}}\) with the following properties.

Each participant's message space admits a measurable bijection with \(\{-1,1\}^d\). The protocol is a one-person-one-message pure-local protocol on the original alphabet \(\mathcal A_d=[d]\times\{-1,1\}\), as specified in \cref{obj:def:tv-private-decision-values}. In particular, its kernels are constant in the history and satisfy
\[
\mathsf K_i(E\mid(j,s),\eta,r)
\le e^\varepsilon
\mathsf K_i(E\mid(j',s'),\eta,r)
\]
for every participant, public seed, history, Borel message event, and pair of original inputs.

The estimator depends on the transcript alone: there is a Borel function \(f\) such that
\[
\widehat D(\mathbf Z^{\mathsf K},R,U)=f(\mathbf Z^{\mathsf K}).
\]
The interval \(J_{\mathrm{TV}}(\mathbf Z^{\mathsf K},R,U)\) is nonempty, connected and closed, has ordered Borel endpoints, and is contained in \([0,1/4]\).

For every \(\theta\in[-1/2,1/2]^d\), under \(n\) iid inputs from \(P_\theta\),
\[
\begin{aligned}
\mathbb E_{P_\theta,\mathsf K}
\left[(\widehat D-D_{\mathrm{TV}}(P_\theta))^2\right]
&\le C_0\,\rho(n\varepsilon^2,d),\\
\Pr_{P_\theta,\mathsf K}
\{D_{\mathrm{TV}}(P_\theta)\in J_{\mathrm{TV}}\}
&\ge0.90,\\
\mathbb E_{P_\theta,\mathsf K}
\left[\operatorname{length}J_{\mathrm{TV}}\right]
&\le C_0\sqrt{\rho(n\varepsilon^2,d)}.
\end{aligned}
\]
Probabilities and expectations include input sampling, channel randomness, the independent public seed, and independent analyst randomization. The same protocol belongs to \(\mathfrak K_{\mathrm{SI}}(n,d,\varepsilon)\).
\end{definitionv}

\begin{propositionv}{prop:signed-causal-bridge}[Signed causal experiment equivalence]
Fix a sample size \(n\in\mathbb N\), a stratum count \(d\in\mathbb N\), and a sampling scheme satisfying:
\begin{itemize}
\item \textbf{(Dimension.)} \(d\ge 2\).
\item \textbf{(Independent participants.)} \cref{obj:ass:iid-people}.
\item \textbf{(Independent randomization.)} \cref{obj:ass:independent-randomness}.
\end{itemize}
Let \(\mathcal V_d\) be the class of causal probability laws in \cref{obj:def:causal-model}. For the stratum \(X\), binary assignment \(A\), observed outcome \(Y\), and potential outcomes \(Y^0,Y^1\), write
\[
S=(2A-1)(2Y-1),\qquad
\mu_{aj}(P)=\mathbb E_P[Y^a\mid X=j],\qquad
\tau_j(P)=\mu_{1j}(P)-\mu_{0j}(P).
\]
Write \(\tau(P)\) for the vector of these contrasts, and define the baseline mean, first-best value, and average absolute contrast by
\[
B(P)=\mathbb E_P[Y],\qquad
V(P)=\frac1d\sum_{j=1}^d\max\{\mu_{0j}(P),\mu_{1j}(P)\},
\qquad
F(\theta)=\frac1d\sum_{j=1}^d|\theta_j|.
\]
The paired signed law \(P_\theta\) has masses
\[
P_\theta(j,s)=\frac{1+s\theta_j}{2d},
\qquad j\in[d],\quad s\in\{-1,1\},
\]
and \(\mathsf U_{2d}\) is the uniform law on this paired alphabet.

For every \(P\in\mathcal V_d\) and every \(j\in[d]\),
\[
\mathbb E_P[S\mid X=j]=\tau_j(P),
\qquad
\mathcal L_P(X,S)=P_{\tau(P)},
\]
and
\[
V(P)
=B(P)+\frac{F(\tau(P))}{2}
=B(P)+\operatorname{TV}(P_{\tau(P)},\mathsf U_{2d}).
\]

For every \(\theta\in[-1/2,1/2]^d\), the symmetric law \(G_\theta\) of \cref{obj:def:symmetric-law} is a probability law in \(\mathcal V_d\), and
\[
V(G_\theta)=\frac12+\frac{F(\theta)}2.
\]
Moreover, for every \(j\in[d]\), \(s\in\{-1,1\}\), and \(a\in\{0,1\}\),
\[
G_\theta(X=j,S=s,A=a)=\frac{P_\theta(j,s)}2,
\]
and
\[
2Y-1=S(2A-1)
\qquad G_\theta\text{-almost surely}.
\]

For every privacy parameter \(\varepsilon\in\mathbb R\), the following two protocol transformations preserve both the sequential and noninteractive classes of \cref{obj:def:sequential-class,obj:def:noninteractive-class}. Let \(\mathcal H_i\) denote the message-history space preceding participant \(i\), and let \(\mathcal R\) denote the declared space of the public-seed random variable \(R\). For every original-record protocol \(Q=(Q_i)_{i=1}^n\), define its paired-alphabet protocol \(\mathsf K\) by averaging over the fair assignment bit:
\[
\mathsf K_i(\,\cdot\mid(j,s),h,r)
=
\frac12\sum_{a\in\{0,1\}}
Q_i\!\left(
\,\cdot\mid
\left(j,a,\frac{1+s(2a-1)}2\right),h,r
\right),
\qquad
h\in\mathcal H_i,\quad r\in\mathcal R.
\]
Sequential membership of \(Q\) at \(\varepsilon\) implies sequential membership of \(\mathsf K\) at \(\varepsilon\), and noninteractive membership of \(Q\) implies noninteractive membership of \(\mathsf K\) at \(\varepsilon\). For every \(\theta\in[-1/2,1/2]^d\), the joint law of the public seed \(R\) and message vector \(\mathbf Z\) generated by \(Q\) under \(G_\theta\) and the fixed sampling scheme equals the joint law generated by \(\mathsf K\) from \(n\) independent \(P_\theta\) inputs and an independent public seed with the same law.

Conversely, for every paired-alphabet protocol \(\mathsf K\), define its original-record protocol \(Q\) by
\[
Q_i(\,\cdot\mid(j,a,y),h,r)
=
\mathsf K_i(\,\cdot\mid(j,(2a-1)(2y-1)),h,r),
\qquad
h\in\mathcal H_i,\quad r\in\mathcal R.
\]
Sequential membership of \(\mathsf K\) at \(\varepsilon\) implies sequential membership of \(Q\) at \(\varepsilon\), and noninteractive membership of \(\mathsf K\) implies noninteractive membership of \(Q\) at \(\varepsilon\). For every \(\theta\in[-1/2,1/2]^d\), the joint public-seed and message law generated by this \(Q\) under \(G_\theta\) and the fixed sampling scheme equals that generated by \(\mathsf K\) from \(n\) independent \(P_\theta\) inputs and an independent public seed with the same law.
\end{propositionv}

\begin{definitionv}{def:vector-kernel}[Sign-vector kernel $Q^{\mathrm{vec}}$]
The sign-vector message kernel is
\[
Q^{\mathrm{vec}}(z\mid O=(j,A,Y))
=
2^{-d}\bigl(1+\delta(2A-1)(2Y-1)z_j\bigr),
\qquad z\in\{-1,1\}^d,
\]
where \(\delta=\tanh(\varepsilon/2)\) is the privacy attenuation. A sampler draws the coordinates independently: every coordinate other than \(j\) is a fair sign, and coordinate \(j\) has mean \(\delta S\), where
\[
S=(2A-1)(2Y-1).
\]
\end{definitionv}

\begin{definitionv}{def:private-moments}[Private moments $\mathsf U_{kj}$]
For a block of \(m\) participants with scaled coordinate messages \(W_{ij}\), define
\[
\mathsf U_{kj}
=
\binom{m}{k}^{-1}
\sum_{\substack{J\subseteq\{1,\ldots,m\}\\ |J|=k}}
\prod_{i\in J}W_{ij}
\quad(1\le k\le m),
\qquad
\mathsf U_{0j}=1.
\]
These quantities are computed by elementary-symmetric recursions.
\end{definitionv}

\begin{algorithmv}{def:frontier-handle}[Public-resource welfare construction $\mathscr H$]
The construction \(\mathscr H\) takes the participant records, sample size \(n\ge2\), dimension \(d\), and privacy parameter \(\varepsilon\) as inputs; its tuning parameters are fixed below, and all auxiliary quantities and summation indices are local to the construction.
\begin{enumerate}
\item Fix the numerical constants, effective private resource, and logarithmic scales:
\[
\eta_*=\frac1{1024},
\qquad L_0=4096,
\qquad t=n\varepsilon^2,
\qquad L=\log(ed),
\qquad z_*=\frac{d^2L}{t},
\qquad w=\log(e+z_*).
\]
Define the squared-error rate function by
\[
\rho(t,d)=
\begin{cases}
\displaystyle\frac{d^2}{tL},&t\ge d^2L,\\[3pt]
\displaystyle\min\left\{1,
\left[\frac{\log(e+d^2L/t)}{L}\right]^2\right\},
&0<t<d^2L.
\end{cases}
\]
Use a singleton public seed and ignore analyst randomization.

\item If \(n=2\), each participant sends a constant message, and set
\[
\widehat V=\frac12.
\]
For this branch, proceed directly to the converse tuning and interval output in the final two steps. For \(n>2\), define the block sizes by
\[
m=\lfloor n/3\rfloor,
\qquad m_0=n-2m.
\]
In the fixed participant order, allocate a baseline block of size \(m_0\), a disjoint sign-vector pilot block of size \(m\), and a disjoint sign-vector evaluation block of size \(m\).

\item For \(n>2\), define the privacy attenuation and scaled-message magnitude by
\[
\delta=\tanh(\varepsilon/2),
\qquad b=\frac d\delta.
\]
Each baseline participant sends a binary message with law
\[
\Pr(\mathsf H_i=h\mid O)
=\frac{1+h\delta(2Y-1)}2,
\qquad h\in\{-1,1\},
\qquad 1\le i\le m_0.
\]
Define the baseline mean estimate by
\[
\widehat B
=\frac12+\frac1{2\delta m_0}
\sum_{i=1}^{m_0}\mathsf H_i.
\]

\item For \(n>2\), each participant in either vector block sends one message with law
\[
Q^{\mathrm{vec}}(z\mid O=(j,A,Y))
=
2^{-d}\bigl(1+\delta(2A-1)(2Y-1)z_j\bigr),
\qquad z\in\{-1,1\}^d,
\]
using independent local randomness. Write \(Z^{\mathrm{pil}}_{ij}\) and \(Z^{\mathrm{ev}}_{ij}\) for coordinate \(j\) of participant \(i\)'s message in the two blocks, and define the scaled-message matrices by
\[
W^{\mathrm{pil}}_{ij}=bZ^{\mathrm{pil}}_{ij},
\qquad
W^{\mathrm{ev}}_{ij}=bZ^{\mathrm{ev}}_{ij},
\qquad 1\le i\le m,\quad 1\le j\le d.
\]
The participant order, binary and sign-vector message alphabets, and kernels are predetermined. This is one fixed noninteractive protocol on the original participant inputs for both protocol classes \(\mathfrak Q_{\mathrm{NI}}\) and \(\mathfrak Q_{\mathrm{SI}}\).

\item For \(n>2\), define the noise scale, polynomial calibration factor, and coordinate means by
\[
\sigma^2=\frac{b^2}{m},
\qquad \sigma=\sqrt{b^2/m},
\qquad H=\log(e+\sigma^2L),
\]
\[
M_j=\frac1m\sum_{i=1}^m W^{\mathrm{pil}}_{ij},
\qquad
\overline W_j=\frac1m\sum_{i=1}^m W^{\mathrm{ev}}_{ij},
\qquad 1\le j\le d.
\]
Define the evaluation-block moments by
\[
\mathsf U_{vj}
=
\binom{m}{v}^{-1}
\sum_{\substack{J\subseteq\{1,\ldots,m\}\\ |J|=v}}
\prod_{i\in J}W^{\mathrm{ev}}_{ij}
\quad(1\le v\le m),
\qquad
\mathsf U_{0j}=1.
\]
Compute these moments by elementary-symmetric recursions.

\item For each even degree \(D\ge2\) selected below, define the Chebyshev polynomials and the absolute-value approximation by
\[
\mathcal T_0(x)=1,
\qquad \mathcal T_1(x)=x,
\qquad
\mathcal T_{v+1}(x)
=2x\mathcal T_v(x)-\mathcal T_{v-1}(x)
\quad(v\ge1),
\]
\[
p_D^{\mathrm{abs}}(x)
=
\frac2\pi+\frac4\pi\sum_{v=1}^{D/2}
\frac{(-1)^{v+1}\mathcal T_{2v}(x)}{4v^2-1}
=
\sum_{v=0}^D c_{D,v}x^v.
\]
For the positive radius \(a_{\mathrm{rad}}\) selected below, define the coordinate polynomial estimate by
\[
\widehat p_{a_{\mathrm{rad}},D,j}
=
\sum_{v=0}^D
c_{D,v}a_{\mathrm{rad}}^{\,1-v}\mathsf U_{vj}.
\]

\item For \(n>2\), define the estimate \(\widehat F\) of the average absolute treatment contrast \(F(\tau(P))\) by taking the first applicable branch below.

If \(L<L_0\), set
\[
\widehat F=\frac1d\sum_{j=1}^d|\overline W_j|.
\]
Otherwise, if \(d^2L\le t\), define
\[
D=2\left\lfloor\frac{L}{2048}\right\rfloor,
\qquad T_*=4\sigma\sqrt L,
\qquad a_{\mathrm{rad}}=8\sigma\sqrt L,
\]
and set
\[
\widehat F
=
\frac1d\sum_{j=1}^d
\left[
\widehat p_{a_{\mathrm{rad}},D,j}
\mathbf1_{\{|M_j|\le T_*\}}
+
\operatorname{sign}(M_j)\overline W_j
\mathbf1_{\{|M_j|>T_*\}}
\right],
\qquad \operatorname{sign}(0)=0.
\]
Otherwise, if \(L/(1024H)<4\), set
\[
\widehat F=\frac14.
\]
Otherwise, define
\[
a_{\mathrm{rad}}=\frac12,
\qquad
D=2\left\lfloor\frac{L}{2048H}\right\rfloor,
\]
and set
\[
\widehat F
=\frac1d\sum_{j=1}^d
\widehat p_{a_{\mathrm{rad}},D,j}.
\]

\item For \(n>2\), combine the baseline and absolute-contrast estimates and clip the resulting welfare estimate:
\[
\widetilde V=\widehat B+\frac12\widehat F,
\qquad
\widehat V=\max\{1/4,\min\{3/4,\widetilde V\}\}.
\]

\item For every allowed resource tuple, define the converse amplitude and integer degrees by
\[
(\alpha,k)=
\begin{cases}
\left(\displaystyle\frac1{100}\sqrt{z_*},
\,2\lceil16L\rceil+1\right),&z_*\le1,\\[3pt]
\left(\displaystyle\frac1{100},
\,2\lceil16L/w\rceil+1\right),&z_*>1,
\end{cases}
\qquad K=k-1.
\]
Use these values in the scaled symmetric moment-matching measures, their independent coordinate product priors, and the symmetric causal laws \(G_\theta\).

\item For every \(n\ge2\), define the interval radius by
\[
q_*=\sqrt{10^{17}\rho(t,d)}.
\]
Output the welfare estimate and connected closed interval
\[
\widehat V=
\begin{cases}
1/2,&n=2,\\
\max\{1/4,\min\{3/4,\widehat B+\widehat F/2\}\},&n>2,
\end{cases}
\qquad
\left[
\max\{1/4,\widehat V-q_*\},
\min\{3/4,\widehat V+q_*\}
\right].
\]
\end{enumerate}
\end{algorithmv}

\begin{theoremv}{thm:uniform-private-value-frontiers}[Uniform private welfare frontiers]*
This result is conditional on the cited, unformalized Bernstein absolute-approximation limit. If \(e_K=\inf_{\deg p\le K}\sup_{|x|\le1}\lvert |x|-p(x)\rvert\) is the best uniform error for approximating \(|x|\) on \([-1,1]\), then \(2k e_{2k}\to\beta_*\) for a constant \(0.278<\beta_*<0.286\). See \citet[Section 3.1]{CaiLow2011Nonsmooth} and \citet{Bernstein1913}.

There are universal constants \(0<c\le C_0<\infty\) with the following properties. For every sample size \(n\), dimension \(d\), privacy budget \(\varepsilon\), and sampling scheme satisfying
\begin{itemize}
\item \textbf{(Resources.)} \(n\ge2\), \(d\ge2\), and \(0<\varepsilon\le1\);
\item \textbf{(Sampling.)} \cref{obj:ass:iid-people};
\item \textbf{(Randomization.)} \cref{obj:ass:independent-randomness};
\end{itemize}
let \(t=n\varepsilon^2\) be the effective private resource and \(L=\log(ed)\) the log-dimension factor. For either protocol class \(\mathcal C\in\{\mathrm{NI},\mathrm{SI}\}\), define
\[
r_{\mathcal C}(n,d,\varepsilon)=\rho(t,d)=
\begin{cases}
\dfrac{d^2}{tL},&t\ge d^2L,\\[4pt]
\displaystyle\min\left\{1,
\left[\frac{\log(e+d^2L/t)}{L}\right]^2\right\},
&0<t<d^2L,
\end{cases}
\qquad
h_{\mathcal C}(n,d,\varepsilon)=\sqrt{\rho(t,d)}.
\]
The minimax squared-error risk \(\mathfrak R_{\mathcal C}\) and globally \(0.90\)-honest connected interval length \(\mathfrak H_{\mathcal C}\), defined in \cref{obj:def:risk,obj:def:honest-length} for this sampling scheme, satisfy
\[
c\,r_{\mathcal C}(n,d,\varepsilon)
\le \mathfrak R_{\mathcal C}(n,d,\varepsilon)
\le C_0\,r_{\mathcal C}(n,d,\varepsilon),
\qquad
c\,h_{\mathcal C}(n,d,\varepsilon)
\le \mathfrak H_{\mathcal C}(n,d,\varepsilon)
\le C_0\,h_{\mathcal C}(n,d,\varepsilon).
\]

For each class, a noninteractive \(\varepsilon\)-private protocol \(Q\), a transcript-only estimator \(T\), and a connected closed interval decision \(I\) with endpoints in \([1/4,3/4]\) attain these upper bounds simultaneously: for every \(P\in\mathcal V_d\), the declared causal model,
\[
\mathbb E_P\!\left[(T(\mathbf Z,R,U)-V(P))^2\right]
\le C_0\,r_{\mathcal C}(n,d,\varepsilon),
\qquad
\mathbb P_P\!\left\{V(P)\in I(\mathbf Z,R,U)\right\}\ge0.90,
\]
\[
\mathbb E_P\!\left[\operatorname{length}
 I(\mathbf Z,R,U)\right]
\le C_0\,h_{\mathcal C}(n,d,\varepsilon).
\]
Here expectations and probabilities use the specified sampling scheme and protocol. The resource-selected explicit noninteractive welfare procedure also satisfies these three guarantees with the same \(C_0\) and the NI rates.

For both classes, risk and squared honest length obey the quantitative comparisons
\[
\frac{c^2}{C_0}\,
\mathfrak R_{\mathcal C}(n,d,\varepsilon)
\le
\mathfrak H_{\mathcal C}(n,d,\varepsilon)^2
\le
\frac{C_0^2}{c}\,
\mathfrak R_{\mathcal C}(n,d,\varepsilon).
\]
The class comparisons are
\[
1\le
\frac{\mathfrak R_{\mathrm{NI}}(n,d,\varepsilon)}
     {\mathfrak R_{\mathrm{SI}}(n,d,\varepsilon)}
\le\frac{C_0}{c},
\qquad
1\le
\frac{\mathfrak H_{\mathrm{NI}}(n,d,\varepsilon)}
     {\mathfrak H_{\mathrm{SI}}(n,d,\varepsilon)}
\le\frac{C_0}{c}.
\]

For every \(d\ge2\) and \(0<t<d^2\log(ed)\), the saturation condition is exactly
\[
\rho(t,d)=1
\quad\Longleftrightarrow\quad
t\le\frac{d^2\log(ed)}{ed-e}.
\]

Along every sequence of allowed resource tuples
\((n_k,d_k,\varepsilon_k)\), write \(t_k=n_k\varepsilon_k^2\). Then
\[
\rho(t_k,d_k)\longrightarrow0
\quad\Longleftrightarrow\quad
t_k\longrightarrow\infty
\ \text{and}\
\frac{\log(1+d_k^2/t_k)}{\log(ed_k)}\longrightarrow0.
\]
Moreover,
\[
\rho(t_k,d_k)\asymp t_k^{-1}
\quad\Longleftrightarrow\quad
\left(\exists a>0:\ t_k\ge a\ \text{eventually}\right)
\ \text{and}\
\left(\exists M\in\mathbb R:\ \min\{t_k,d_k\}\le M\
\text{eventually}\right),
\]
where \(\asymp\) denotes eventual comparability by positive finite constants. In particular, when \(t_k\to\infty\),
\[
\rho(t_k,d_k)\asymp t_k^{-1}
\quad\Longleftrightarrow\quad
\exists M\in\mathbb N:\ d_k\le M\ \text{eventually}.
\]
For every family of sampling schemes satisfying
\cref{obj:ass:iid-people,obj:ass:independent-randomness} at each \(n,d\), and for either fixed class \(\mathcal C\), these consistency and private-parametric characterizations also hold for the corresponding minimax risk. Honest length tends to zero under the same consistency condition, and its private-parametric order is \(t_k^{-1/2}\) under the same comparability conditions.

Finally, for every \(0<a\le b<\infty\), there are constants \(0<c\le C_0<\infty\) and an integer \(D_0\), possibly depending on \(a,b\), such that for every \(d\ge\max\{D_0,2\}\) and every \(t\) with \(a\le t/d^2\le b\),
\[
c\left[\frac{\log\log(ed)}{\log(ed)}\right]^2
\le \rho(t,d)\le
C_0\left[\frac{\log\log(ed)}{\log(ed)}\right]^2.
\]
Consequently, along allowed sequences with \(d_k\to\infty\) and \(t_k/d_k^2\) in a fixed positive finite interval, either class has
\[
\mathfrak R_{\mathcal C}(n_k,d_k,\varepsilon_k)
\asymp
\left[\frac{\log\log(ed_k)}{\log(ed_k)}\right]^2,
\qquad
\mathfrak H_{\mathcal C}(n_k,d_k,\varepsilon_k)
\asymp
\frac{\log\log(ed_k)}{\log(ed_k)}.
\]
\end{theoremv}
% CAUSALSMITH-CITED-SCOPE-BEGIN thm:uniform-private-value-frontiers
\verificationfootnotetext{\textbf{Formalization scope.} This result uses the cited conclusion \citep{CaiLow2011Nonsmooth} (Section 3.1, definition of the Bernstein constant and immediately following paragraph, page 8). The cited source proof is not formalized here: Lean verifies its use and all remaining steps. If that cited conclusion is false, the portions of this result that depend on it are not certified.}
% CAUSALSMITH-CITED-SCOPE-END thm:uniform-private-value-frontiers

\begin{theoremv}{thm:two-cell-calibration}[Two-cell risk and length calibration]
There exist universal constants \(0<c\le C_0<\infty\) such that, for every sample size \(n\), privacy budget \(\varepsilon\), and sampling scheme for two-cell observed records satisfying
\begin{itemize}
\item \textbf{(Resources.)} \(n\ge2\) and \(0<\varepsilon\le1\).
\item \textbf{(Sampling.)} \cref{obj:ass:iid-people}.
\item \textbf{(Randomness.)} \cref{obj:ass:independent-randomness}.
\end{itemize}
the minimax squared-error risk and globally \(0.90\)-honest expected interval length defined in \cref{obj:def:risk,obj:def:honest-length}, evaluated under this sampling scheme, satisfy, for every \(\mathcal C\in\{\mathrm{NI},\mathrm{SI}\}\),
\[
c\min\{1,(n\varepsilon^2)^{-1}\}
\le \mathfrak R_{\mathcal C}(n,2,\varepsilon)
\le C_0\min\{1,(n\varepsilon^2)^{-1}\},
\]
\[
c\min\{1,(\sqrt n\,\varepsilon)^{-1}\}
\le \mathfrak H_{\mathcal C}(n,2,\varepsilon)
\le C_0\min\{1,(\sqrt n\,\varepsilon)^{-1}\}.
\]

Moreover, for each such sampling scheme, \(n\), and \(\varepsilon\), there exist a total noninteractive protocol \(Q\in\mathfrak Q_{\mathrm{NI}}(n,2,\varepsilon)\) on the original observed records, an estimator \(T(\mathbf Z,R,U)\) equal to a Borel function of the message transcript \(\mathbf Z\) alone, and a connected closed interval \(I(\mathbf Z,R,U)\) with ordered Borel endpoints in \([1/4,3/4]\), such that, simultaneously for every \(P\in\mathcal V_2\), the causal model in \cref{obj:def:causal-model},
\[
\mathbb E\!\left[(T(\mathbf Z,R,U)-V(P))^2\right]
\le C_0\min\{1,(n\varepsilon^2)^{-1}\},
\]
\[
\mathbb P\!\left\{V(P)\in I(\mathbf Z,R,U)\right\}\ge0.90,
\qquad
\mathbb E\!\left[\operatorname{length}(I(\mathbf Z,R,U))\right]
\le C_0\min\{1,(\sqrt n\,\varepsilon)^{-1}\}.
\]
The expectations and probability are taken under the decision law induced by the sampling scheme, \(Q\), and \(P\).
\end{theoremv}

\begin{definitionv}{def:tv-corollary-models}[Total-variation models $\mathcal P_d^{\mathrm{pair}}$ and $\Delta_{2d}$]
Define the finite paired alphabet and its probability simplex by
\[
\mathcal A_d=[d]\times\{-1,1\},
\qquad
\Delta_{2d}
=
\left\{
\mathsf p:\mathcal A_d\to[0,1]:
\sum_{j=1}^d\sum_{s\in\{-1,1\}}\mathsf p(j,s)=1
\right\}.
\]
Here \(\mathsf p\) is a probability law on the paired alphabet. Define the paired family and its uniform-reference total-variation functional by
\[
\mathcal P_d^{\mathrm{pair}}
=
\{P_\theta:\theta\in[-1/2,1/2]^d\}
\subseteq\Delta_{2d},
\]
\[
D_{\mathrm{TV}}(\mathsf p)
=
\frac12\sum_{j=1}^d\sum_{s\in\{-1,1\}}
\left|\mathsf p(j,s)-\frac1{2d}\right|.
\]
The reference law \(\mathsf U_{2d}\) is public and uniform on \(\mathcal A_d\). In either statistical model
\[
\mathcal M=\mathcal P_d^{\mathrm{pair}}
\quad\text{or}\quad
\mathcal M=\Delta_{2d},
\]
the data are \(n\) independent, identically distributed draws from one unknown law \(\mathsf p\in\mathcal M\), and the estimand is \(D_{\mathrm{TV}}(\mathsf p)\). The causal model \(\mathcal V_d\) retains its full-data laws and welfare estimand \(V(P)\).
\end{definitionv}

\begin{definitionv}{def:tv-private-decision-values}[Private total-variation decision values $\mathfrak R^{\mathrm{TV}}_{\mathcal C}$ and $\mathfrak H^{\mathrm{TV}}_{\mathcal C}$]
The class \(\mathfrak K_{\mathcal C}(n,d,\varepsilon)\), for \(\mathcal C\in\{\mathrm{NI},\mathrm{SI}\}\), consists of one-message-per-participant protocols on the paired input alphabet \(\mathcal A_d=[d]\times\{-1,1\}\). A protocol \(\mathsf K\) satisfies the following conditions:
\begin{itemize}
\item \textbf{(Spaces and kernels.)} The protocol has an arbitrary standard-Borel public-seed space with probability law \(\lambda\), standard-Borel participant message spaces, and participant-specific message kernels \(\mathsf K_i\) conditional on the input, preceding message history, and public seed.
\item \textbf{(Local privacy.)} For every participant \(i\), public seed \(r\), preceding message history \(\eta\), Borel message event \(E\), \(j,j'\in[d]\), and \(s,s'\in\{-1,1\}\),
\[
\mathsf K_i(E\mid(j,s),\eta,r)
\le e^\varepsilon\mathsf K_i(E\mid(j',s'),\eta,r).
\]
\item \textbf{(Public specification.)} All protocol choices depend only on \(n,d,\varepsilon\), the specified model, and public bounds.
\item \textbf{(Interaction.)} For \(\mathcal C=\mathrm{NI}\), every kernel is constant in its message history; participant-specific kernels and shared public seeds are permitted. For \(\mathcal C=\mathrm{SI}\), kernels may depend on preceding messages, and each participant is accessed once.
\item \textbf{(Independent randomizers.)} The public seed \(R\) and analyst randomizer \(U\) have laws
\[
R\sim\lambda,
\qquad
U\sim\operatorname{Uniform}[0,1],
\]
and are mutually independent and independent of the independent, identically distributed inputs.
\end{itemize}

Write
\[
\mathbf Z^{\mathsf K}=(Z_i^{\mathsf K})_{i=1}^n
\]
for the resulting message vector. An estimate \(\widehat D\) is any Borel real-valued function of \((\mathbf Z^{\mathsf K},R,U)\). An interval \(J_{\mathrm{TV}}\) is any nonempty connected closed interval contained in \([0,1]\), with ordered Borel endpoints that depend only on the same information.

For either statistical model
\[
\mathcal M\in\{\mathcal P_d^{\mathrm{pair}},\Delta_{2d}\},
\]
the minimax squared-error risk and globally honest expected interval length for the uniform-reference total-variation distance \(D_{\mathrm{TV}}(\mathsf p)\) are
\[
\begin{aligned}
\mathfrak R^{\mathrm{TV}}_{\mathcal C}(\mathcal M;n,d,\varepsilon)
&=
\inf_{\mathsf K\in\mathfrak K_{\mathcal C}(n,d,\varepsilon)}
\inf_{\widehat D}
\sup_{\mathsf p\in\mathcal M}
\mathbb E_{\mathsf p,\mathsf K}
\bigl[(\widehat D-D_{\mathrm{TV}}(\mathsf p))^2\bigr],\\
\mathfrak H^{\mathrm{TV}}_{\mathcal C}(\mathcal M;n,d,\varepsilon)
&=
\inf_{\mathsf K\in\mathfrak K_{\mathcal C}(n,d,\varepsilon)}
\inf_{\substack{J_{\mathrm{TV}}:\,
\inf_{\mathsf p\in\mathcal M}
\Pr_{\mathsf p,\mathsf K}
\{D_{\mathrm{TV}}(\mathsf p)\in J_{\mathrm{TV}}\}\ge0.90}}
\sup_{\mathsf p\in\mathcal M}
\mathbb E_{\mathsf p,\mathsf K}
\operatorname{length}J_{\mathrm{TV}}.
\end{aligned}
\]
Here \(\operatorname{length}J_{\mathrm{TV}}\) is the difference between its ordered endpoints. Probabilities and expectations include input sampling, channel randomness, public randomness, and analyst randomization. Singleton intervals are permitted. For the paired model, the attaining intervals specified below are contained in \([0,1/4]\).
\end{definitionv}

\begin{theoremv}{thm:paired-uniform-tv-frontier}[Paired-uniform total-variation frontier]*
This result is conditional on the cited, unformalized Bernstein absolute-approximation limit. If \(e_K=\inf_{\deg p\le K}\sup_{|x|\le1}\lvert |x|-p(x)\rvert\) is the best uniform error for approximating \(|x|\) on \([-1,1]\), then \(2k e_{2k}\to\beta_*\) for a constant \(0.278<\beta_*<0.286\). See \citet[Section 3.1]{CaiLow2011Nonsmooth} and \citet{Bernstein1913}.

For the paired family in \cref{obj:def:tv-corollary-models}, every integer \(d\ge2\) and every \(\theta\in[-1/2,1/2]^d\) satisfy
\[
 P_\theta(j,s)=\frac{1+s\theta_j}{2d},
 \qquad
 D_{\mathrm{TV}}(P_\theta)
 =\operatorname{TV}(P_\theta,\mathsf U_{2d})
 =\frac12F(\theta)
 =\frac{1}{2d}\sum_{j=1}^{d}|\theta_j|.
\]
Here \(\mathsf U_{2d}\) is the uniform law on the paired alphabet and \(F(\theta)\) is the average absolute contrast.

There are universal real constants \(0<c\le C_0<\infty\) such that the following conclusions hold for every resource tuple satisfying
\begin{itemize}
\item \textbf{(Sample size.)} \(n\) is an integer with \(n\ge2\).
\item \textbf{(Dimension.)} \(d\) is an integer with \(d\ge2\).
\item \textbf{(Privacy.)} \(0<\varepsilon\le1\).
\end{itemize}
For either \(\mathcal C\in\{\mathrm{NI},\mathrm{SI}\}\), let
\[
 t=n\varepsilon^2,\qquad
 r_{\mathcal C}(n,d,\varepsilon)=\rho(t,d),\qquad
 h_{\mathcal C}(n,d,\varepsilon)=\sqrt{\rho(t,d)},
\]
where \(\rho\) is the evaluated rate defined in \cref{obj:def:frontier-handle}. The decision values of \cref{obj:def:tv-private-decision-values} obey
\[
\begin{aligned}
 c\,r_{\mathcal C}(n,d,\varepsilon)
 &\le
 \mathfrak R^{\mathrm{TV}}_{\mathcal C}
   (\mathcal P_d^{\mathrm{pair}};n,d,\varepsilon)
 \le C_0\,r_{\mathcal C}(n,d,\varepsilon),\\
 c\,h_{\mathcal C}(n,d,\varepsilon)
 &\le
 \mathfrak H^{\mathrm{TV}}_{\mathcal C}
   (\mathcal P_d^{\mathrm{pair}};n,d,\varepsilon)
 \le C_0\,h_{\mathcal C}(n,d,\varepsilon).
\end{aligned}
\]

For each such resource tuple there exists a pure \(\varepsilon\)-locally private noninteractive paired-input protocol \(\mathsf K\), with each participant's message space in bijection with \(\{-1,1\}^d\), together with a transcript-only estimator \(\widehat D\) and a connected closed interval \(J_{\mathrm{TV}}\subseteq[0,1/4]\), such that, for every \(\mathsf p\in\mathcal P_d^{\mathrm{pair}}\), under independent sampling from \(\mathsf p\),
\[
\begin{aligned}
 \mathbb E_{\mathsf p,\mathsf K}
   [(\widehat D-D_{\mathrm{TV}}(\mathsf p))^2]
 &\le C_0\,r_{\mathcal C}(n,d,\varepsilon),\\
 \mathbb P_{\mathsf p,\mathsf K}
   \{D_{\mathrm{TV}}(\mathsf p)\in J_{\mathrm{TV}}\}
 &\ge0.90,\\
 \mathbb E_{\mathsf p,\mathsf K}
   [\operatorname{length}(J_{\mathrm{TV}})]
 &\le C_0\,h_{\mathcal C}(n,d,\varepsilon).
\end{aligned}
\]
These guarantees also hold for the concrete resource-selected sign-vector polynomial protocol and its total-variation estimator and interval, with the construction specified in \cref{obj:def:frontier-handle,obj:thm:uniform-private-value-frontiers}: this protocol is noninteractive and pure \(\varepsilon\)-locally private, and its squared risk, coverage, and expected length satisfy the three displayed bounds with \(\mathcal C=\mathrm{NI}\).

For the same universal constants and every admissible resource tuple,
\[
\begin{aligned}
1
&\le
\frac{\mathfrak R^{\mathrm{TV}}_{\mathrm{NI}}
  (\mathcal P_d^{\mathrm{pair}};n,d,\varepsilon)}
 {\mathfrak R^{\mathrm{TV}}_{\mathrm{SI}}
  (\mathcal P_d^{\mathrm{pair}};n,d,\varepsilon)}
\le \frac{C_0}{c},\\
1
&\le
\frac{\mathfrak H^{\mathrm{TV}}_{\mathrm{NI}}
  (\mathcal P_d^{\mathrm{pair}};n,d,\varepsilon)}
 {\mathfrak H^{\mathrm{TV}}_{\mathrm{SI}}
  (\mathcal P_d^{\mathrm{pair}};n,d,\varepsilon)}
\le \frac{C_0}{c}.
\end{aligned}
\]

The rate has the following exact resource properties. Write \(L=\log(ed)\). For every integer \(d\ge2\) and every \(0<t<d^2L\),
\[
 \rho(t,d)=1
 \quad\Longleftrightarrow\quad
 t\le\frac{d^2L}{ed-e}.
\]
For every sequence of admissible resource tuples
\((n_k,d_k,\varepsilon_k)\), write
\[
 t_k=n_k\varepsilon_k^2,\qquad L_k=\log(ed_k).
\]
Then
\[
 \rho(t_k,d_k)\longrightarrow0
 \quad\Longleftrightarrow\quad
 t_k\longrightarrow\infty
 \ \text{and}\
 \frac{\log(1+d_k^2/t_k)}{L_k}\longrightarrow0.
\]
Moreover,
\[
 \rho(t_k,d_k)\asymp t_k^{-1}
 \quad\Longleftrightarrow\quad
 \bigl(\exists a>0:\ t_k\ge a\ \text{eventually}\bigr)
 \ \text{and}\
 \bigl(\exists M\in\mathbb R:\ \min\{t_k,d_k\}\le M\
       \text{eventually}\bigr).
\]
Here \(\asymp\) means two-sided comparison by fixed positive finite constants for all sufficiently large \(k\). If \(t_k\to\infty\), this equivalence becomes
\[
 \rho(t_k,d_k)\asymp t_k^{-1}
 \quad\Longleftrightarrow\quad
 \exists M\in\mathbb N:\ d_k\le M\ \text{eventually}.
\]
For every pair of real numbers \(0<a\le b\), there exist real constants \(0<c\le C_0<\infty\) and an integer \(D_0\) such that, for every integer \(d\ge D_0\) with \(d\ge2\) and every real \(t\) satisfying \(a\le t/d^2\le b\),
\[
 c\left(\frac{\log L}{L}\right)^2
 \le \rho(t,d)
 \le C_0\left(\frac{\log L}{L}\right)^2.
\]

Finally, these sequence properties transfer to the actual minimax values for each fixed \(\mathcal C\in\{\mathrm{NI},\mathrm{SI}\}\). Along every admissible sequence,
\[
\begin{aligned}
 \mathfrak R^{\mathrm{TV}}_{\mathcal C}
   (\mathcal P_{d_k}^{\mathrm{pair}};n_k,d_k,\varepsilon_k)
   \longrightarrow0
 &\quad\Longleftrightarrow\quad
 t_k\longrightarrow\infty
 \ \text{and}\\
 \frac{\log(1+d_k^2/t_k)}{L_k}\longrightarrow0,\\
 \mathfrak H^{\mathrm{TV}}_{\mathcal C}
   (\mathcal P_{d_k}^{\mathrm{pair}};n_k,d_k,\varepsilon_k)
   \longrightarrow0
 &\quad\Longleftrightarrow\quad
 t_k\longrightarrow\infty
 \ \text{and}\\
 \frac{\log(1+d_k^2/t_k)}{L_k}\longrightarrow0.
\end{aligned}
\]
The squared-error value satisfies
\[
 \mathfrak R^{\mathrm{TV}}_{\mathcal C}
   (\mathcal P_{d_k}^{\mathrm{pair}};n_k,d_k,\varepsilon_k)
   \asymp t_k^{-1}
 \quad\Longleftrightarrow\quad
 \bigl(\exists a>0:\ t_k\ge a\ \text{eventually}\bigr)
 \ \text{and}\
 \bigl(\exists M\in\mathbb R:\ \min\{t_k,d_k\}\le M\
       \text{eventually}\bigr).
\]
Under \(t_k\to\infty\), equivalently,
\[
 \mathfrak R^{\mathrm{TV}}_{\mathcal C}
   (\mathcal P_{d_k}^{\mathrm{pair}};n_k,d_k,\varepsilon_k)
   \asymp t_k^{-1}
 \quad\Longleftrightarrow\quad
 \exists M\in\mathbb N:\ d_k\le M\ \text{eventually}.
\]
For each fixed \(\mathcal C\) and every \(0<a\le b\), there exist real constants \(0<c\le C_0<\infty\) and an integer \(D_0\) such that every admissible tuple with \(d\ge D_0\) and \(a\le n\varepsilon^2/d^2\le b\) satisfies
\[
\begin{aligned}
 c\left(\frac{\log L}{L}\right)^2
 &\le
 \mathfrak R^{\mathrm{TV}}_{\mathcal C}
   (\mathcal P_d^{\mathrm{pair}};n,d,\varepsilon)
 \le C_0\left(\frac{\log L}{L}\right)^2,\\
 \sqrt c\left|\frac{\log L}{L}\right|
 &\le
 \mathfrak H^{\mathrm{TV}}_{\mathcal C}
   (\mathcal P_d^{\mathrm{pair}};n,d,\varepsilon)
 \le \sqrt{C_0}\left|\frac{\log L}{L}\right|.
\end{aligned}
\]
\end{theoremv}
% CAUSALSMITH-CITED-SCOPE-BEGIN thm:paired-uniform-tv-frontier
\verificationfootnotetext{\textbf{Formalization scope.} This result uses the cited conclusion \citep{CaiLow2011Nonsmooth} (Section 3.1, definition of the Bernstein constant and immediately following paragraph, page 8). The cited source proof is not formalized here: Lean verifies its use and all remaining steps. If that cited conclusion is false, the portions of this result that depend on it are not certified.}
% CAUSALSMITH-CITED-SCOPE-END thm:paired-uniform-tv-frontier

\begin{theoremv}{thm:full-simplex-tv-converse}[Full-simplex total-variation lower bounds]*
This result is conditional on the cited, unformalized Bernstein absolute-approximation limit. If \(e_K=\inf_{\deg p\le K}\sup_{|x|\le1}\lvert |x|-p(x)\rvert\) is the best uniform error for approximating \(|x|\) on \([-1,1]\), then \(2k e_{2k}\to\beta_*\) for a constant \(0.278<\beta_*<0.286\). See \citet[Section 3.1]{CaiLow2011Nonsmooth} and \citet{Bernstein1913}.

Let \(\Delta_{2d}\) be the probability simplex on the paired alphabet \(\mathcal A_d=[d]\times\{-1,1\}\), and let
\[
D_{\mathrm{TV}}(\mathsf p)=\operatorname{TV}(\mathsf p,\mathsf U_{2d}),
\]
where \(\mathsf U_{2d}\) is the uniform reference law. There is a universal constant \(c>0\) such that, for every sample size \(n\), dimension \(d\), and privacy budget \(\varepsilon\) satisfying
\begin{itemize}
\item \textbf{(Sample size.)} \(n\ge2\).
\item \textbf{(Dimension.)} \(d\ge2\).
\item \textbf{(Privacy budget.)} \(0<\varepsilon\le1\).
\end{itemize}
and every protocol-class label \(\mathcal C\in\{\mathrm{NI},\mathrm{SI}\}\), the private decision values of \cref{obj:def:tv-private-decision-values}, with iid inputs drawn from \(\mathsf p\in\Delta_{2d}\), satisfy
\[
\begin{aligned}
\mathfrak R^{\mathrm{TV}}_{\mathcal C}(\Delta_{2d};n,d,\varepsilon)
&\ge c\,r_{\mathcal C}(n,d,\varepsilon),\\
\mathfrak H^{\mathrm{TV}}_{\mathcal C}(\Delta_{2d};n,d,\varepsilon)
&\ge c\,h_{\mathcal C}(n,d,\varepsilon),
\end{aligned}
\]
where the squared-error and connected-length rates are
\[
r_{\mathcal C}(n,d,\varepsilon)=\rho(n\varepsilon^2,d),
\qquad
h_{\mathcal C}(n,d,\varepsilon)=\sqrt{r_{\mathcal C}(n,d,\varepsilon)}.
\]
The interval value uses connected closed intervals in \([0,1]\) with coverage at least \(0.90\) at every law in \(\Delta_{2d}\). The protocol classes permit arbitrary public seeds, participant-specific kernels, standard-Borel message spaces, and, for the sequential class, dependence on preceding messages, with one message per participant and pure local privacy.
\end{theoremv}
% CAUSALSMITH-CITED-SCOPE-BEGIN thm:full-simplex-tv-converse
\verificationfootnotetext{\textbf{Formalization scope.} This result uses the cited conclusion \citep{CaiLow2011Nonsmooth} (Section 3.1, definition of the Bernstein constant and immediately following paragraph, page 8). The cited source proof is not formalized here: Lean verifies its use and all remaining steps. If that cited conclusion is false, the portions of this result that depend on it are not certified.}
% CAUSALSMITH-CITED-SCOPE-END thm:full-simplex-tv-converse

\begin{theoremv}{lem:bounded-mean-concentration}[Bounded mean concentration]
Let \(m\ge1\) be an integer, let \(b\in\mathbb R\), and let \(H_1,\ldots,H_m\) be mutually independent, measurable real random variables on a common probability space. Suppose that
\[
-b\le H_i\le b
\]
at every sample point, for each \(i=1,\ldots,m\). Then, for every real \(x\ge0\),
\[
\begin{aligned}
\Pr\left\{m^{-1}\sum_{i=1}^m(H_i-\mathbb E H_i)\ge x\right\}
&\le
\begin{cases}
\exp\!\left(-\dfrac{mx^2}{2b^2}\right),& b\ne0,\\
1,& b=0,
\end{cases}\\
\Pr\left\{m^{-1}\sum_{i=1}^m(H_i-\mathbb E H_i)\le -x\right\}
&\le
\begin{cases}
\exp\!\left(-\dfrac{mx^2}{2b^2}\right),& b\ne0,\\
1,& b=0.
\end{cases}
\end{aligned}
\]
\end{theoremv}

\begin{theoremv}{lem:cai-low-chebyshev-approximation}[Chebyshev approximation of absolute value]
Let \(D\ge 2\) be an even integer. Define the Chebyshev polynomials by
\[
\mathcal T_0(x)=1,\qquad \mathcal T_1(x)=x,\qquad
\mathcal T_{v+1}(x)=2x\mathcal T_v(x)-\mathcal T_{v-1}(x)
\quad (v\ge 1),
\]
and define the polynomial \(q_D\) and its monomial coefficients \(a_{Dv}\) by
\[
q_D(x)=\frac{2}{\pi}
+\frac{4}{\pi}\sum_{v=1}^{D/2}
\frac{(-1)^{v+1}\mathcal T_{2v}(x)}{4v^2-1}
=\sum_{v=0}^{D}a_{Dv}x^v.
\]
Then, for every \(x\in[-1,1]\),
\[
\bigl|q_D(x)-|x|\bigr|\le \frac{2}{\pi(D+1)},
\]
and, for every integer \(v\) with \(0\le v\le D\),
\[
|a_{Dv}|\le 2^{3D/2}.
\]
\end{theoremv}

\begin{theoremv}{lem:cai-low-moment-duality}[Symmetric moment matching duality]
For every positive even integer \(K\), define the best uniform polynomial approximation error
\[
e_K=\inf_{\deg p\le K}\sup_{|x|\le 1}\bigl|p(x)-|x|\bigr|,
\]
where the infimum is over real polynomials \(p\). There exist probability measures \(\xi_0,\xi_1\) on \(\mathbb R\), each supported on \([-1,1]\) and invariant under reflection \(x\mapsto -x\), such that
\[
\int x^\ell\,d\xi_1(x)=\int x^\ell\,d\xi_0(x)
\qquad\text{for every integer }0\le \ell\le K,
\]
and
\[
\int |x|\,d\xi_1(x)-\int |x|\,d\xi_0(x)=2e_K.
\]
\end{theoremv}

\begin{lemmav}{lem:private-moment-and-dependence}[Private moments and column dependence]
Let \(n,d,m\) be nonnegative integers and let \(\varepsilon\) be the privacy budget. Suppose that:
\begin{itemize}
\item \textbf{(Sampling.)} The sampling scheme satisfies \cref{obj:ass:iid-people}.
\item \textbf{(Population.)} The full-data law \(P\) belongs to \(\mathcal V_d\), as defined in \cref{obj:def:causal-model}.
\item \textbf{(Privacy calibration.)} \(0<\varepsilon\le1\).
\item \textbf{(Dimension and block size.)} \(d\ge2\) and \(0<m\le n\).
\end{itemize}
Apply the vector channel \(Q^{\mathrm{vec}}\) once per participant, with a singleton public seed. Define the privacy attenuation and scaled-message magnitude by
\[
\delta=\tanh(\varepsilon/2),
\qquad
b=\frac d\delta.
\]
For an \(m\)-participant message block, write \(Z_{ij}\in\{-1,1\}\) for coordinate \(j\) of participant \(i\)'s sign-vector message and define the scaled messages by
\[
W_{ij}=bZ_{ij},
\qquad 1\le i\le m,\quad j\in[d].
\]
Write \(\tau_j(P)\) for the cellwise treatment contrast,
\[
\tau_j(P)=\mu_{1j}(P)-\mu_{0j}(P).
\]
All expectations, variances, and covariances below are taken under the block message law induced by \(P\) and \(Q^{\mathrm{vec}}\).

The protocol satisfies
\[
Q^{\mathrm{vec}}\in\mathfrak Q_{\mathrm{NI}}(n,d,\varepsilon),
\]
where membership has the full-record privacy and noninteraction requirements of \cref{obj:def:noninteractive-class}. For every \(1\le i\le m\) and \(j\in[d]\),
\[
\mathbb E[W_{ij}]=\tau_j(P),
\qquad
\mathbb E[W_{ij}^{2}]=b^{2}.
\]
For every \(1\le i\le m\) and distinct \(j,j'\in[d]\),
\[
\mathbb E[W_{ij}W_{ij'}]=0.
\]
The private moments \(\mathsf U_{kj}\) of \(W\), defined in \cref{obj:def:private-moments}, satisfy, for every integer \(0\le k\le m\) and every \(j\in[d]\),
\[
\mathbb E[\mathsf U_{kj}]=\tau_j(P)^k.
\]
For every nonnegative integer \(k\) with \(2k\le m\) and every \(j\in[d]\),
\[
\mathbb E[\mathsf U_{kj}^{2}]
\le
\left(\tau_j(P)^2+\frac{2kb^2}{m}\right)^k.
\]

For any deterministic measurable column statistics \(T_j^{\mathrm{col}}\), each a function of \((W_{1j},\ldots,W_{mj})\), every pair of distinct coordinates \(j,j'\in[d]\) satisfies
\[
\left|\operatorname{Cov}
\left(T_j^{\mathrm{col}},T_{j'}^{\mathrm{col}}\right)\right|
\le
\frac{|\tau_j(P)\tau_{j'}(P)|}
{\sqrt{(b^2-\tau_j(P)^2)(b^2-\tau_{j'}(P)^2)}}
\sqrt{\operatorname{Var}(T_j^{\mathrm{col}})
      \operatorname{Var}(T_{j'}^{\mathrm{col}})}.
\]
The same collection satisfies
\[
\operatorname{Var}\left(\frac1d\sum_{j=1}^{d}T_j^{\mathrm{col}}\right)
\le
\frac2d\max_{j\in[d]}\operatorname{Var}(T_j^{\mathrm{col}}).
\]
\end{lemmav}

\begin{lemmav}{lem:finite-private-polynomial-calibration}[Finite private polynomial calibration]
Let \(n,d\) be sample size and dimension, \(\varepsilon\) the privacy budget, and \(m\) the evaluation-block size. Suppose that:
\begin{itemize}
\item \textbf{(Resources.)} \(n\ge2\), \(d\ge2\), \(0<\varepsilon\le1\), and \(0<m\le n\).
\item \textbf{(Sampling.)} The sampling scheme satisfies \cref{obj:ass:iid-people}.
\item \textbf{(Causal model.)} The full-data law \(P\) belongs to \(\mathcal V_d\), as defined in \cref{obj:def:causal-model}.
\item \textbf{(Bounded mean concentration.)} For every positive integer \(m\), every \(b>0\), and every collection of independent measurable real random variables \(H_1,\ldots,H_m\) taking values in \([-b,b]\), the following bounds hold for every \(x\ge0\):
\[
 \Pr\!\left\{\frac1m\sum_{i=1}^m(H_i-\mathbb EH_i)\ge x\right\}
 \le \exp\!\left(-\frac{mx^2}{2b^2}\right),
 \qquad
 \Pr\!\left\{\frac1m\sum_{i=1}^m(H_i-\mathbb EH_i)\le -x\right\}
 \le \exp\!\left(-\frac{mx^2}{2b^2}\right).
\]
\item \textbf{(Chebyshev approximation.)} For every even integer \(D\ge2\), the explicit Chebyshev absolute-value approximation, with monomial coefficients \(a_{Dv}\),
\[
 q_D(x)=\sum_{v=0}^D a_{Dv}x^v,
\]
satisfies
\[
 \sup_{x\in[-1,1]}|q_D(x)-|x||\le\frac{2}{\pi(D+1)},
 \qquad
 |a_{Dv}|\le 2^{3D/2}\quad(0\le v\le D).
\]
\end{itemize}

Set the privacy attenuation, scaled-message magnitude, logarithmic dimension, calibration constant, and block noise scale to
\[
 \delta=\tanh(\varepsilon/2),\qquad
 b=\frac d\delta,\qquad
 L=\log(ed),\qquad
 \eta_*=\frac1{1024},\qquad
 \sigma^2=\frac{b^2}{m},\qquad
 \sigma=\sqrt{\sigma^2}.
\]
Use independent participant releases from the signed-vector kernel
\[
 S=(2A-1)(2Y-1),\qquad
 Q^{\mathrm{vec}}(z\mid O=(j,A,Y))
 =2^{-d}(1+\delta S z_j),
 \qquad z\in\{-1,1\}^d,
\]
and let \(W_{ij}\) be coordinate \(j\) of release \(i\), multiplied by \(b\). On an evaluation block of size \(m\), define the distinct-person moments and polynomial estimates by
\[
 \mathsf U_{0j}=1,\qquad
 \mathsf U_{vj}
 =\binom mv^{-1}
   \sum_{\substack{I\subseteq\{1,\ldots,m\}\\|I|=v}}
      \prod_{i\in I}W_{ij}
 \quad(1\le v\le m),
\]
\[
 p_{a,D}(x)=a q_D(x/a),\qquad
 \widehat p_{a,D,j}
 =\sum_{v=0}^D a_{Dv}a^{1-v}\mathsf U_{vj}.
\]
The target coordinate and its average absolute magnitude are
\[
 \tau_j(P)=\mu_{1j}(P)-\mu_{0j}(P),\qquad
 F(\tau(P))=\frac1d\sum_{j=1}^d|\tau_j(P)|.
\]

For every even integer \(D\ge2\) satisfying
\[
 2D\le m,\qquad
 H=\log(e+\sigma^2L),\qquad
 D\le\frac{\eta_*L}{H},
\]
take \(a=1/2\). Then, for every \(j\in[d]\),
\[
 \left|\mathbb E\widehat p_{1/2,D,j}-|\tau_j(P)|\right|
 \le\frac1{D+1},
 \qquad
 \operatorname{Var}\!\left(\frac1d\sum_{j=1}^d
                    \widehat p_{1/2,D,j}\right)
 \le\frac{\exp(9DH)}{2d}.
\]

For every integer \(D\ge0\) satisfying
\[
 L\ge4096,\qquad
 D=2\left\lfloor\frac{\eta_*L}{2}\right\rfloor,\qquad
 2D\le m,\qquad
 2m\le n,
\]
use two independent signed-vector release blocks of size \(m\), with scaled messages \(W^{\mathrm{pil}}\) and \(W^{\mathrm{ev}}\). Define
\[
 M_j=\frac1m\sum_{i=1}^mW^{\mathrm{pil}}_{ij},\qquad
 \overline W_j=\frac1m\sum_{i=1}^mW^{\mathrm{ev}}_{ij},\qquad
 T_*=4\sigma\sqrt L,\qquad a=2T_*,
\]
and form \(\widehat p_{a,D,j}\) from the evaluation-block moments. With \(\operatorname{sign}(0)=0\), put
\[
 H_j=\widehat p_{a,D,j}\mathbf1_{\{|M_j|\le T_*\}}
       +\operatorname{sign}(M_j)\overline W_j
          \mathbf1_{\{|M_j|>T_*\}}.
\]
Then, for every \(j\in[d]\),
\[
 \left|\mathbb EH_j-|\tau_j(P)|\right|
 \le\frac{20000\sigma}{\sqrt L},
 \qquad
 \operatorname{Var}(H_j)
 \le400\sigma^2L\exp(L/100),
\]
and
\[
 \mathbb E\!\left[
   \left(\frac1d\sum_{j=1}^dH_j-F(\tau(P))\right)^2
 \right]
 \le500000000\,\frac{\sigma^2}{L}.
\]
\end{lemmav}

\begin{theoremv}{lem:measurable-kernel-radon-nikodym}[Jointly measurable Radon–Nikodym density]
Let \(H\) and \(Z\) be measurable spaces, and let \(K,L:H\to Z\) be kernels satisfying:
\begin{itemize}
\item \textbf{(Standard Borel target.)} \(Z\) is a standard-Borel space.
\item \textbf{(Finiteness.)} Both \(K\) and \(L\) are finite kernels.
\item \textbf{(Pointwise domination.)} \(K(h,\cdot)\ll L(h,\cdot)\) for every \(h\in H\).
\end{itemize}
Then there exists a jointly measurable function \(f:H\times Z\to[0,\infty]\) such that, for every \(h\in H\) and every measurable set \(E\subseteq Z\),
\[
K(h,E)=\int_E f(h,z)\,L(h,dz).
\]
\end{theoremv}

\begin{lemmav}{lem:adaptive-moment-contraction}[Adaptive higher-order transcript contraction]
Let \(n,d\) be integers and let \(\varepsilon\) be a privacy budget. Suppose that:
\begin{itemize}
\item \textbf{(Parameter regime.)} \(n\ge2\), \(d\ge2\), and \(0<\varepsilon\le1\).
\item \textbf{(Sampling and randomness.)} The sampling scheme satisfies \cref{obj:ass:iid-people,obj:ass:independent-randomness}.
\item \textbf{(Sequential privacy.)} \(Q\in\mathfrak Q_{\mathrm{SI}}(n,d,\varepsilon)\), as specified in \cref{obj:def:sequential-class}.
\item \textbf{(Measurable kernel densities.)} For every measurable parameter space \(H\), standard-Borel output space \(Z\), and finite kernels \(K,L\) from \(H\) to \(Z\) satisfying \(K(h,\cdot)\ll L(h,\cdot)\) for every \(h\), there exists a jointly measurable \(f:H\times Z\to[0,\infty]\) such that, for every \(h\) and measurable \(E\subseteq Z\),
\[
K(h,E)=\int_E f(h,z)\,L(h,dz).
\]
\end{itemize}
Set the privacy contraction scale to
\[
\beta=\frac{e^\varepsilon-1}{2d}.
\]
There exists a jointly measurable real-valued density version \(p^Q(\theta,r)(\mathbf z)\), defined for all \(\theta\in\mathbb R^d\), public seeds \(r\), and transcripts \(\mathbf z\), with the following properties. For every \(\theta\in[-1/2,1/2]^d\) and every \(r\), it is nonnegative at every transcript and gives the conditional transcript law under \(G_\theta\) relative to the zero-contrast reference:
\[
\Pr_Q(d\mathbf z\mid r,G_\theta)
=
p^Q(\theta,r)(\mathbf z)\,
\Pr_Q(d\mathbf z\mid r,G_0).
\]
For every \(\theta,r,\mathbf z,j\), its dependence on coordinate \(\theta_j\), with the other coordinates fixed, is a polynomial on \(\mathbb R\) of degree at most \(n\). For every \(\theta\in[-1/2,1/2]^d\), every \(r,j\), and every integer \(1\le k\le n\),
\[
\int
\left|\partial_{\theta_j}^{\,k}p^Q(\theta,r)(\mathbf z)\right|
\,\Pr_Q(d\mathbf z\mid r,G_0)
\le
k!\sqrt{\binom nk}\,\beta^k.
\]
For every \(\theta\in\mathbb R^d\), every \(r,\mathbf z,j\), and every integer \(k>n\),
\[
\partial_{\theta_j}^{\,k}p^Q(\theta,r)(\mathbf z)=0.
\]

For every amplitude \(0<\alpha\le1/2\), let \(\nu_0,\nu_1\) be probability laws supported on \([-\alpha,\alpha]\), fixed independently of the public seed. Define their seed-transcript mixtures by
\[
\mathsf M_\ell^Q(dr,d\mathbf z)
=
\lambda(dr)
\int_{[-\alpha,\alpha]^d}
\Pr_Q(d\mathbf z\mid r,G_\theta)\,
\nu_\ell^{\otimes d}(d\theta),
\qquad \ell\in\{0,1\}.
\]
For every integer \(k\ge1\) such that
\[
\int u^q\,\nu_0(du)=\int u^q\,\nu_1(du),
\qquad q=0,\ldots,k-1,
\]
the following conclusions hold:
\[
\operatorname{TV}(\mathsf M_0^Q,\mathsf M_1^Q)
\le
\min\left\{1,\,
d\alpha^k\sqrt{\binom nk}\,\beta^k\right\},
\qquad 1\le k\le n,
\]
and
\[
\mathsf M_0^Q=\mathsf M_1^Q,
\qquad k>n.
\]
Moreover, every Borel measurable seed-transcript decision \(f(r,\mathbf z)\in[0,1]\), interpreted as a randomized test's acceptance probability, satisfies
\[
\left|
\int f(r,\mathbf z)\,\mathsf M_0^Q(dr,d\mathbf z)
-
\int f(r,\mathbf z)\,\mathsf M_1^Q(dr,d\mathbf z)
\right|
\le
\operatorname{TV}(\mathsf M_0^Q,\mathsf M_1^Q).
\]
\end{lemmav}

\begin{lemmav}{lem:separated-value-mixtures}[Separated welfare mixture bounds]
Fix a sample size \(n\), a dimension \(d\), a sampling scheme satisfying \cref{obj:ass:iid-people,obj:ass:independent-randomness}, and a law-independent one-message sequential protocol \(Q\) with probability message kernels and standard-Borel seed and message spaces. Write \(\mathcal V_d\) for the full-data probability laws satisfying the declared causal model of uniform strata, fair assignment, consistency, and interior arm means. The welfare target is
\[
V(P)=\frac{1}{d}\sum_{j=1}^{d}
\max\{\mu_{0j}(P),\mu_{1j}(P)\},
\]
where \(\mu_{aj}(P)\) is the conditional mean of potential outcome \(Y^a\) in stratum \(j\).

Consider two probability priors on full-data laws, called the first and second priors. Let \(v_0<v_1\) be their respective welfare centers, and define their separation by
\[
\Delta=v_1-v_0.
\]
Let \(\eta_0,\omega\in\mathbb R\). Suppose:
\begin{itemize}
\item \textbf{(Causal support.)} Each prior assigns probability one to \(\mathcal V_d\).
\item \textbf{(Decision mixtures.)} Mixing the induced joint law of \((\mathbf Z,R,U)\) over either prior yields a probability measure. Let \(N_0\) and \(N_1\) be the respective joint seed-transcript marginals, namely the mixture laws of \((R,\mathbf Z)\).
\item \textbf{(Target concentration.)} Under the first and second priors, respectively,
\[
\Pr\{|V(P)-v_0|>\Delta/8\}\le\eta_0,
\qquad
\Pr\{|V(P)-v_1|>\Delta/8\}\le\eta_0.
\]
\item \textbf{(Transcript proximity.)} The joint seed-transcript mixtures satisfy
\[
\operatorname{TV}(N_0,N_1)\le\omega.
\]
\end{itemize}

Then every Borel real-valued estimator \(T(\mathbf Z,R,U)\) satisfies
\[
\sup_{P\in\mathcal V_d}
\mathbb E_{P,Q}\bigl[(T(\mathbf Z,R,U)-V(P))^2\bigr]
\ge
\frac{9\Delta^2}{128}(1-\omega-2\eta_0)_+,
\]
with squared-error expectations interpreted in the extended nonnegative reals.

Moreover, every connected closed interval \(I(\mathbf Z,R,U)\) with ordered Borel endpoints in \([1/4,3/4]\) and global coverage
\[
\Pr_{P,Q}\{V(P)\in I(\mathbf Z,R,U)\}\ge0.90
\qquad\text{for every }P\in\mathcal V_d
\]
satisfies
\[
\sup_{P\in\mathcal V_d}
\mathbb E_{P,Q}\operatorname{length}I(\mathbf Z,R,U)
\ge
\frac{3\Delta}{4}(0.80-2\eta_0-\omega)_+.
\]
Here \(\operatorname{length}I\) is the upper endpoint minus the lower endpoint, and \(x_+=\max\{x,0\}\).
\end{lemmav}

\begin{lemmav}{lem:dimension-free-private-two-point}[Dimension-free private two-point lower bounds]
Assume the following measurable Radon--Nikodym property: for every measurable parameter space \(H\), every standard-Borel output space \(Z\), and every pair of finite kernels \(K,L\) from \(H\) to \(Z\) satisfying \(K(h,\cdot)\ll L(h,\cdot)\) for every \(h\in H\), there exists a jointly measurable density \(f:H\times Z\to[0,\infty]\) such that, for every \(h\in H\) and every measurable \(E\subseteq Z\),
\[
K(h,E)=\int_E f(h,z)\,L(h,dz).
\]
Then there exists a universal constant \(c_1>0\) such that, for every sample size \(n\in\mathbb N\), dimension \(d\in\mathbb N\), privacy budget \(\varepsilon\in\mathbb R\), and sampling scheme satisfying
\begin{itemize}
\item \textbf{(Resources.)} \(n\ge2\), \(d\ge2\), and \(0<\varepsilon\le1\).
\item \textbf{(Participant sampling.)} \cref{obj:ass:iid-people}.
\item \textbf{(Randomization.)} \cref{obj:ass:independent-randomness}.
\end{itemize}
the minimax squared-error risk \(\mathfrak R_{\mathcal C}(n,d,\varepsilon)\) and globally \(0.90\)-honest connected interval length \(\mathfrak H_{\mathcal C}(n,d,\varepsilon)\), defined in \cref{obj:def:risk,obj:def:honest-length} for that sampling scheme, satisfy, for every \(\mathcal C\in\{\mathrm{NI},\mathrm{SI}\}\),
\[
\mathfrak R_{\mathcal C}(n,d,\varepsilon)
\ge c_1\min\left\{1,\frac{1}{n\varepsilon^2}\right\},
\qquad
\mathfrak H_{\mathcal C}(n,d,\varepsilon)
\ge c_1\min\left\{1,\frac{1}{\sqrt n\,\varepsilon}\right\}.
\]
\end{lemmav}
